%
% grundlagen.tex -- Basics
%
% (c) 2020 Prof Dr Andreas Müller, Hochschule Rapperswil
%
% !TEX root = ../../buch.tex
% !TEX encoding = UTF-8
%

\section{Grundlagen\label{qa:section:einleitung}}
\kopfrechts{Grundlagen}
\subsubsection{Addition und Multiplikation komplexer Zahlen}
Sei eine positive reele Zahl $x$: Wenn $x$ mit $i$ multipliziert 
wird, dann entsteht aus einer rein reeler Zahl eine rein imaginäre 
Zahl $ix$. Dies kann als eine 90° Drehung im Gegenuhrzeigersinn 
interpretiert werden. Wird nun das Ergebnis $ix$ ein weiteres Mal 
mit $i$ multipliziert, sodass $iix = i^2x$, dann greift die oben 
genannte Regel und $-x$ entsteht. Die Zahl wurde somit um weitere 
90° im Gegenuhrzeigersinn gedreht. Die selbe Regel gilt auch 
allgemein: Selbst wenn $x$ keine reele positive Zahl, sondern eine 
komplexe Zahl ist, geschieht die Drehung auf gleicher Weise. 

Für eine ganze 360° Drehung muss $x$ vier mal mit $i$ multipliziert 
werden. Potenzen mit Basis $i$ wiederholen sich also nach vier 
Iterationen.

% Tabelle und Bild einfügen

Mit Multiplikationen können nicht nur Drehungen vollzogen werden. 
Statt eine Zahl mit $i$ zu multiplizieren, kann mit einer 
Multiplikation von zwei komplexen Zahlen $c_1, c_2$ eine präzise 
Platzierung mit bestimmter Drehung und Skalierung durchgeführt 
werden. Dies geschieht distributiv, wobei das Ergebnis wieder in 
Real- und Imaginärteil aufgeteilt wird:
\begin{equation*}
    (a + ib) \cdot (c + id) 
    = ac + iad + ibc - bd = 
    (ac - bd) + i(ad + bc)
\end{equation*}

Die Addition zweier komplexer Zahlen gleicht dem Vorgang einer 
Vektoraddition, wobei die Zahlen komponentenweise zusammengezählt   
werden:
\begin{equation*}
    (a + ib) + (d + id) = a+c + ib + id = (a+c) + i(b+d)
\end{equation*}
Auch hier wird das Ergebnis auf den Real- und Imaginärteil 
aufgeteilt. Quaternionen verhalten sich sehr ähnlich, wie im 
nächsten Abschnitt beschrieben ist.

\subsubsection{Addition und Multiplikationen von Quaternionen}
Für Quaternionen werden zwei neue Regeln eingeführt. Die erste 
Regel ist die selbe wie bei den Komplexen Zahlen. Sie besagt, dass 
das Quadrat einer imaginären Achse jeweils $-1$ ergibt:
\begin{equation*}
    i^2 = j^2 = k^2 = -1
\end{equation*}
Die zweite Regel führt eine Abhängigkeit zwischen den imaginären 
Achsen ein. Eine Folge der drei Imaginärachsen in alphabetischer 
Reihenfolge $ijk$ ergibt ebenfalls $-1$. Dies gilt auch für 
zyklische Verschiebungen, wie in
\begin{equation*}
    ijk = kij = jki = -1
\end{equation*}
beschrieben. Damit entstehen weitere Eigenschaften: Da $ijk = kk = 
-1$, muss $ij = k$ entsprechen. Dasselbe gilt auch für $ki = j$ und 
$jk = i$. Die Reihenfolge einer Multiplikation spielt dabei eine 
grosse Rolle, denn $ik \neq ki$. Dies kann wie folgt nachgewiesen 
werden, indem das folgende Experiment durchgeführt wird:
\begin{align*}
    ik &= i \cdot 1 \cdot k \\
    ik &= i \cdot -(j)^2 \cdot k \\
    ik &= -i \cdot j^2 \cdot k \\
    ik &= -\underbrace{i \cdot j}_k \cdot 
    \underbrace{j \cdot k}_{i} \\
    ik  &= -ki\\
\end{align*}
Wenn die Reihenfolge der Multiplikation vertauscht wird, dann 
ändert sich das Vorzeichen. Quaternionen sind also nicht 
kommutatiiv. Dies kann intuitiv verstanden werden, wenn man sich 
eine Drehung um die $x$-Achse und eine Drehung um die $y$-Achse 
vorstellt. Die Reihenfolge der beiden Drehungen spielt eine Rolle, 
da das Ergebnis unterschiedlich ausfällt.

Die Addition von zwei Quaternionen funktioniert analog zu den 
komplexen Zahlen. Die Komponenten werden komponentenweise addiert:
\begin{equation*}
    (a + ib + jc + kd) + (e + if + jg + kh) 
    = (a+e) + i(b+f) + j(c+g) + k(d+h)
\end{equation*}
Auch das Multiplizieren zweier Quaternionen funktioniert analog zu 
den komplexen Zahlen. Die Komponenten werden distributiv 
multipliziert und das Ergebnis wieder in Real- und Imaginärteile 
aufgeteilt:
\begin{align*}
    (a + ib + jc + kd) \cdot (e + if + jg + kh) 
    &= ae + aif + ajg + akh \\
    &+ ibe + i^2bf + ijbg + ikbh \\
    &+ jce + jif + j^2cg + jkch \\
    &+ kde + kif + kjcg + k^2dh \\
    &= (ae - bf - cg - dh) \\
    &+ i(af + be + ch - dg) \\
    &+ j(ag - bh + ce + df) \\
    &+ k(ah + bg - cf + de)
\end{align*}

\subsubsection{Konjugation, Norm und Inverse}

Auch für Quaternionen existiert die Komplex-Konjugation. Wie bei 
den komplexen Zahlen wird der Realteil nicht verändert, während die 
Imaginärteile negiert werden:
\begin{equation*}
    q^* = (a + ib + jc + kd)^* = a - ib - jc - kd
\end{equation*}
Die L2-Norm eines Quaternions ist die Summe der Quadrate aller 
Komponenten. Sie wird wie folgt berechnet:
\begin{equation*}
    |q|^2 = \|a + ib + jc + kd\|^2 = a^2 + b^2 + c^2 + d^2
\end{equation*}
Die Norm lässt sich auch mit der Konjugation berechnen:
\begin{equation*}
    |q|^2 = q \cdot q^* 
    = (a + ib + jc + kd) \cdot (a - ib - jc - kd)
\end{equation*}
Ein Einheitsquaternion ist ein Quaternion mit der Norm 1. Sie 
bilden die Einheitskugel im vierdimensionalen Raum, wie der 
Einheitskreis im komplexen Raum. In den folgenden Abschnitten wird 
gezeigt, dass vor allem Einheitsquaternionen von grosser 
Wichtigkeit sind.

Um die Inverse eines Quaternions zu berechnen, werden die Norm und 
die Konjugation benötigt. Das Ziel ist, dass $q \cdot q^{-1} = 1$ 
ergibt. Diese Gleichung kann wie folgt umgestellt werden:
\begin{align*}
    q \cdot q^{-1} &= 1 \\
    q^{-1} &= \frac{1}{q} \\
    q^{-1} &= \frac{q^*}{q \cdot q^*} \\
    q^{-1} &= \frac{q^*}{\|q\|^2}
\end{align*}
Was auffällt ist, dass die Inverse eines Einheitsquaternions gleich 
der Konjugation ist. Damit lassen sich Einheitsquaternionen sehr 
einfach invertieren.